%% ============================================================================
\section{Conclusions}
\label{ch:conclusions}
%% ============================================================================

This thesis asked whether a profiler-adaptive \code{DataLoader} that combines
bottleneck classification, tar-shard staging and predictive prefetching can improve
GPU utilisation and tile throughput for large-scale satellite-image segmentation on
HPC systems.

\subsection{Summary of contributions}

\paragraph{The workload is I/O-bound, and this is measurable.}
Section~\ref{ch:materials} characterised the segmentation corpus not by its class
balance - the conventional description - but by its file layout and access
pattern.
The Bing 1k configuration comprises 5{,}025 image and 5{,}025 mask files stored
individually on a shared parallel filesystem, giving a per-sample cost dominated by
metadata operations rather than by bytes transferred.
Baseline profiling showed a CPU-to-GPU time ratio across all three architectures
large enough that the data path is a plausible constraint (8--20$\times$ under
the corrected profiler methodology, Table~\ref{tab:rq1-breakdown}), and showed
that changing two loader flags moves GPU utilisation by 24--51 percentage points
depending on architecture (Table~\ref{tab:rq1-loader-util}) -- a substantially
larger effect than the IoU movement itself, which was mixed in direction across
architectures (DeepLabV3+ +3.2pp, MA-Net +2.6pp, SegFormer $-$2.5pp, default to
tuned) and smaller in magnitude than the reference study's architectural search.
The loader configuration is where this workload's headroom actually is, but that
headroom shows up in throughput, not directly in accuracy.

\paragraph{\pads{}.}
Section~\ref{ch:methodology} presented a drop-in loader that profiles its own
behaviour over a bounded warm-up window, classifies the dominant bottleneck by an
inspectable rule set over measured timings, selects among loose, shard and stage
policies, and derives its prefetch depth from measured staging and consumption times
subject to a scratch-capacity constraint.
The design rests on a property of supervised training that is usually left unused:
once the epoch permutation is drawn, future file accesses are known exactly, so
prefetching is exact rather than speculative.

\paragraph{An evaluation that isolates mechanisms.}
Sections~\ref{ch:results} and~\ref{ch:ablation} separate two claims a single
headline number would conflate. The staging primitive itself works: sharding
alone beats the loose baseline on GPU utilisation for every architecture tested
(Table~\ref{tab:ablation-components}), and a full pre-stage ceiling reaches the
best GPU utilisation of any method tested for two of three architectures
(Table~\ref{tab:rq3-endtoend}). \pads{}'s adaptive decision layer, however, does
not reliably capture that benefit: it selected \code{loose} in every RQ3
method-5 run and in 10 of 12 Section~\ref{ch:ablation} threshold-sweep runs.
Traced to its cause (Section~\ref{sec:res-profiler-bias}), this is not evidence
that adaptive staging cannot work -- it is evidence that the specific warm-up
probe measuring the decision was biased by roughly an order of magnitude, and a
validated fix (discard 5 warm-up batches before measuring) is already
implemented. The honest summary is: the primitive works; the instrument that
decided when to use it was measurably wrong, and now is not.

\paragraph{A reproduction, and the defects it surfaced.}
Appendix~\ref{app:reproduction} reproduced the reference configuration and
established the accuracy floor against which any pipeline change is judged.
Appendix~\ref{app:defects} documents the defects encountered, two of which are
consequential beyond this thesis: the staging policy was applied to loader
construction but not to dataset construction, so the selected policy never reached
the data path; and no archive existed in a layout the shard reader could open.
Together these mean that staging results collected before these fixes describe
policy \emph{decisions} rather than executions.
That finding is reported here in full because it invalidates prior measurements
rather than merely improving on them.

\subsection{Answers to the research questions}

\begin{description}
  \item[RQ1 - When does the loader become the bottleneck?]
    Always, in this configuration: \Tdata{} exceeds \Tgpu{} by 8--20$\times$
    across all three architectures under the corrected measurement
    (Table~\ref{tab:rq1-breakdown}). The predicted ordering by \Tgpu{} -- that
    the lightest architecture (SegFormer) should hit the bottleneck hardest --
    held only partially: it is confirmed for SegFormer specifically, but
    DeepLabV3+ and MA-Net's default-configuration GPU utilisation
    (Table~\ref{tab:rq1-loader-util}) were not ordered as the heavier-model
    prediction anticipated.
  \item[RQ2 - Do tar shards reduce many-small-file overhead?]
    Mechanistically, yes: \code{loose} issues roughly two filesystem operations
    per sample (Table~\ref{tab:rq2-fsops}), consistent with RQ1's finding that
    \Tdata{} is metadata-bound rather than bandwidth-bound. This translates into
    wall-clock benefit when the sharded/staged reader is actually invoked
    (Table~\ref{tab:ablation-components}), which -- see RQ3 below -- is not as
    often as the adaptive design intends.
  \item[RQ3 - Does predictive staging outperform fixed streaming?]
    Not reliably, as measured: \pads{} matched or exceeded the full-prestage
    ceiling for only one of three architectures (MA-Net), and its adaptive
    profiler selected \code{loose} in every method-5 run in
    Table~\ref{tab:rq3-endtoend}. Section~\ref{sec:res-profiler-bias} traces
    this to a specific, quantified, and now-corrected measurement defect in the
    profiler rather than a limitation of the staging mechanism itself, which
    Table~\ref{tab:ablation-components} shows does help when actually invoked.
  \item[RQ4 - Can prefetch depth be auto-tuned?]
    No, not by the present formula: every auto-tuned depth landed at
    $d\approx170$--$235$ (Table~\ref{tab:rq4-auto}), one to two orders of
    magnitude outside the $1$--$8$ range the fixed sweep tests and where a fixed
    depth was shown to help. The derivation is not merely imprecise; it depends
    on a quantity, $T_{stage}$, shown to vary by three to four orders of
    magnitude between repeated measurements on identical hardware.
\end{description}

\subsection{What this work does not claim}

Three limits are worth stating plainly, because the value of the contribution
depends on where its boundary lies.

First, the result is specific to a corpus of many small files on a contended shared
filesystem.
A workload whose data already resides on fast node-local storage, or whose samples
are large enough that per-file overhead amortises, has no bottleneck for \pads{} to
remove; the loader is designed to detect exactly this and degrade to the ordinary
\code{DataLoader}, and reporting that it does so is part of the claim rather than an
admission against it.

Second, \pads{} optimises the data path and does not improve the model.
Segmentation quality is treated throughout as a constraint to be preserved rather
than an objective to be advanced, and any IoU movement beyond the reference standard
deviation is treated as a regression to be explained.

Third, the accuracy figures used for the policy comparisons in
Sections~\ref{ch:results} and~\ref{ch:ablation} are not this workload's ceiling.
A memory-constrained workstation reproduction under an otherwise faithful
configuration attained substantially lower IoU than the reference figure, traced
to the batch-size reduction forced by accelerator memory and its effect on
batch-normalisation statistics (Appendix~\ref{app:reproduction}). That diagnosis
was subsequently confirmed directly, not only by analogy: a full-protocol
replica run on the cluster at the reference's own batch size, epoch budget and
final-epoch weighting reached $72.8$--$73.5\%$ IoU across all three
architectures, within $1$--$1.4$ points of the published $74.17 \pm 0.38$
(Section~\ref{sec:repro-closure}). The policy-comparison figures elsewhere in
this thesis are lower still than either of these, for two further, identified
reasons specific to the main pipeline (an IoU-averaging convention and a
flip-augmentation defect, Appendix~\ref{app:defects}) rather than because the
workload's accuracy ceiling is actually 0.35. The workstation environment
remains adequate for instrumentation and defect analysis, which is the use made
of it in Appendix~\ref{app:defects}, but accuracy claims throughout this thesis
should be read against the closed-gap figure above, not the main pipeline's
uncorrected numbers.

\subsection{Closing remarks}

The broader observation this work supports is methodological, though it is more
precisely stated than an earlier draft of this chapter had it. Effort in deep
learning for remote sensing is conventionally directed at the model, and the
reference study for this workload found that architectural choice moved IoU by
only a few percentage points - plausibly within the noise of random
augmentation. Two loader flags, in the measurements collected this session, do
not clearly move IoU more than that: the effect was mixed in direction and
comparable in magnitude across architectures (+3.2, +2.6 and $-2.5$ percentage
points). What two loader flags move decisively is GPU utilisation - by 24 to 51
percentage points, an effect an order of magnitude larger than anything observed
on accuracy in this thesis, in either direction. The corrected claim is not that
loader configuration is a hidden lever on accuracy; it is that the accelerator's
idle time, not the model, is where this pipeline's headroom actually sits, and
that headroom is measurable in throughput terms regardless of what it does or
does not do to IoU. The adaptive layer meant to capture that headroom
automatically remains the weaker part of this work
(Section~\ref{sec:res-profiler-bias}), and closing that gap between a mechanism
shown to work and a decision procedure shown not to invoke it reliably is the
most direct next step (Section~\ref{ch:futurework}).
